% ==========================================================
% Section: Methodology
% ==========================================================

\section{Methodology}
\label{sec:methodology}

In this section, we formulate the decoupled template framework and its mathematical underpinning for content reusability.

\subsection{Mathematical Formulation}
Let $\mathcal{C}$ represent the set of core scientific content sections, and let $\mathcal{T} = \{T_{\text{custom}}, T_{\text{ieee}}, T_{\text{elsevier}}, T_{\text{acm}}, T_{\text{springer}}\}$ denote the set of target publisher templates. A compiled manuscript $\mathcal{M}_i$ for a given publisher template $T_i \in \mathcal{T}$ is obtained through a deterministic mapping $\Phi$:
\begin{equation}
    \mathcal{M}_i = \Phi(T_i, \mathcal{C}, \mathcal{B}, \mathcal{F})
    \label{eq:manuscript_mapping}
\end{equation}
where $\mathcal{B}$ denotes the bibliography database and $\mathcal{F}$ represents the visual figure assets.

\subsection{Architecture Overview}
The system enforces strict invariance on $\mathcal{C}$:
\begin{equation}
    \frac{\partial \mathcal{C}}{\partial T_i} = 0, \quad \forall T_i \in \mathcal{T}
\end{equation}
This invariance ensures that modifying typography or publisher requirements never introduces discrepancies or side effects into the scientific discourse.
